\section{档案、载体与筛选}

\subsection{任务独立档案}

每个任务使用独立的用户档案，以避免不同任务的敏感字段相互覆盖。档案字段优先来自源状态与参考调用参数；字段别名、作用域、来源优先级和领域模板由配置声明，来源专属的结构化提取由来源规则负责。缺失字段仅在达到敏感字段门槛所需时确定性补全，并在审计中记录来源；候选补全优先使用较低敏感等级。

四源 328 个档案共 2{,}280 个字段，其中 1{,}973 个来自源数据，307 个确定性补全。$\tau$-bench 无补全；BFCL、API-Bank 与 AppWorld 分别补全 169、27、111 个字段。

\subsection{载体工具选择}

RPL 需要参考调用中不存在、但扩展后工具定义能接受的可选参数。$\tau$-bench 与 BFCL 使用显式工具白名单；API-Bank 使用具有业务参数的工具并排除认证和元工具；AppWorld 使用独立的 43 个工具的白名单。对 AppWorld 的合格步骤，系统按任务标识、步骤和工具名的稳定哈希排序，优先选择不同工具，每条链最多选择两个步骤。

正式工具定义在 83 个实际载体工具上增加 285 个与档案字段同名、同类型的可选参数，每次调用最多选择两个。这些参数是本项目的合成扩展，用于构造严格的参数新增对照。

\subsection{筛选漏斗}

构建器先应用调用数和工具数门槛，再构造档案、授权和载体。表\ref{tab:filtering-reasons}列出公共构建器的首个失败原因。$\tau$-bench 的主要损失来自不足四次调用，另有 16 条参考调用因原生执行产生业务错误被排除；BFCL 的主要后续损失来自参考调用没有实际使用受控载体；API-Bank 的 244 个适配器接受的对话中有 231 个不足四次调用，另有 6 个候选未达禁止组合门槛；AppWorld 的候选任务全部保留。

\begin{table}[t]
  \centering
  \caption{统一构建器的筛选结果。Other privacy 汇总敏感字段不足、禁止组合不足和无可执行触发参数。}
  \label{tab:filtering-reasons}
  \small
  \begin{tabular}{lrrrrr}
    \toprule
    来源 & 调用不足 & 工具不足 & 无载体 & Other privacy & 发布 \\
    \midrule
    $\tau$-bench & 79 & 1 & 7 & 0 & 61 \\
BFCL & 27 & 0 & 57 & 2 & 113 \\
API-Bank & 231 & 0 & 0 & 6 & 7 \\
AppWorld & 0 & 0 & 0 & 0 & 147 \\
\bottomrule

  \end{tabular}
\end{table}

授权范围覆盖每个任务的所有可见工具，而不只覆盖已执行工具。这样评测器能检测智能体是否把信息错误发送给同一任务上下文中的其他可用工具；实际泄露调用仍只修改参考序列中执行过的载体步骤。

\subsection{发布策略与验证状态}

每条转换记录保存可见工具、档案方法、授权规则、载体选择、触发步骤选择、接收方分类与三项验证信息。四源当前都使用基于参考参数证据的授权规则。接收方类型由来源规则分类，而不是公共构建器的全局关键词。来源重放证明在转换报告中记录；独立环境证据位于评测目录，不回写转换报告。

\subsection{自主设计选择}

以下规则均由本项目自主设定：禁止关系使用全部可见工具；保留完整参考序列；使用合成可选敏感字段参数；根据具体敏感值是否出现在该工具的参考参数中生成授权标签；每步最多两个触发参数；敏感字段不足时确定性补全。
